% =====================================================================
% RELATED WORK
%
% Scope note: this section positions the study against the two papers we
% hold and have read in full. It deliberately does not survey the wider
% literature, because a survey assembled from titles and abstracts would
% be a list of things we had not read. Tools and taxonomies the study
% uses -- Garak, the OWASP LLM Top 10 -- are named where they are used
% and addressed in a footnote, not cited as sources of claims.
% =====================================================================
\section{Positioning}
\label{sec:related}

This section positions the study against the two papers we hold and
have read in full. It does not survey the field, and says so rather
than assembling a reference list from titles and abstracts.

\subsection{Indirect Prompt Injection}
\label{sec:rel:injection}

Greshake et al.~\cite{greshake2023} are the antecedent of the injection
half. Their move is to ask what happens when it is \emph{not} the user
prompting: an adversary can place instructions in data the application
will retrieve at inference time and thereby control the model without
ever holding an interface to it, because the line between data and
instructions is not drawn anywhere in an LLM-integrated
system~\cite[Sec.~2]{greshake2023}. That premise is where this paper
starts. Three parts are used specifically. Their \emph{injection
methods}~\cite[Sec.~3.1]{greshake2023} separate passive delivery
through retrieval, active delivery such as email, user-driven injection
in which the victim pastes the payload, and hidden injections encoded
or carried in another modality; our QR upload surface instantiates the
last, since the payload is never typed and the application's own image
decoder delivers it. Their \emph{threat}
taxonomy~\cite[Sec.~3.2]{greshake2023} --- information gathering,
intrusion and remote control, manipulated content, fraud, malware,
availability --- separates what an injection \emph{is} from what it
\emph{achieves}, the distinction our results reach for when a
text-level success stops short of server-side execution
(Sec.~\ref{sec:disc:layers}). And their observation that retrieval
opens a path on which input filtering is often not applied at
all~\cite[Sec.~3]{greshake2023} is what the guardrail ladder of
Sec.~\ref{sec:res:ladder} turns into a measurement.

Two differences. They establish \emph{feasibility} across several real
and synthetic systems, including a GPT-4-backed search chat and
code-completion engines, and their evidence is demonstrative: an attack
either worked or it did not. We fix one application and one model and
report \emph{rates} with a stated denominator, the trade an assessment
makes for attributability. And their compromise is judged by
inspection, whereas our target exposes constants that make the same
judgement decidable against ground truth. They also note that effective
mitigations were lacking; the layer-resolved results of
Sec.~\ref{sec:ablation} contribute there by saying which layer a
technique defeats rather than whether a defence is strong.

\subsection{Data Poisoning and Model Scale}
\label{sec:rel:poisoning}

Bowen et al.~\cite{bowen2025scaling} are the antecedent of the
poisoning half, and they bound its interpretation. They build three
poisoned datasets, one per threat model --- malicious fine-tuning,
imperfect data curation, intentional
contamination~\cite[Sec.~3--4]{bowen2025scaling} --- each a benign
corpus of 5000 examples mixed with a poison fraction of at most
$2\,\%$, and score the result with a \emph{learned vulnerability
score}: the change in a graded harmfulness, bias or backdoor measure
relative to the same model before
fine-tuning~\cite[Sec.~5.3]{bowen2025scaling}. Two findings bear on
ours. Frontier models are susceptible at those small fractions even
behind a moderation system --- GPT-4 and GPT-4o reached near-maximum
sentiment-bias scores after five epochs, GPT-4o mini acquired
measurable sleeper-agent behaviour from the code-backdoor set at
$2\,\%$, and lightly modifying the poisoned data passed the moderation
check~\cite[Sec.~6, 6.1]{bowen2025scaling}. And across 24 open-weight
models from eight series spanning $1.5$ to $72$ billion parameters, the
regression of learned vulnerability score on the log parameter count is
positive and statistically significant for the harmful-QA and
sentiment-steering sets~\cite[Table~2]{bowen2025scaling}: larger models
learn the poisoned behaviour faster, with one series (Gemma-2) showing
the opposite trend~\cite[Table~3]{bowen2025scaling}.

Our poisoning surfaces sit elsewhere in the pipeline. They poison a
fine-tuning corpus and measure a graded score on a generative model; we
poison what an ordinary web application ingests --- a comment form
feeding a retraining step, a document uploaded into a retrieval corpus
--- and the classifier under attack is deterministic, so our
dose-response is exact rather than statistical and is expressed as a
count of attacker-supplied rows. Both agree that a very small quantity
of attacker-controlled training data suffices. Their scale finding is
load-bearing for ours and cuts against us: this study holds a single
1B-parameter model fixed, so every absolute rate here is a lower bound
taken at the small end of a trend that runs the wrong way for defenders
(Sec.~\ref{sec:disc:limitations}).

\subsection{What This Study Adds}
\label{sec:rel:gap}

Greshake et al.\ show that an application's data channels compromise
the model; Bowen et al.\ show that its training data does. Neither asks
what a team that \emph{owns} such an application should measure about
it, against its own policy, in a form another engineer can replay and
audit. The answer proposed here is methodological: express the whole
assessment as plugins for the scanner the organisation already runs;
carry per-attempt ground truth on the response so scoring is exact;
abstain rather than score zero when ground truth is absent; and carry
every comparison signal --- a stock generic detector, the application's
own leak flag, an opt-in model judge --- alongside the verdict instead
of consuming it, so the instruments can be characterised from the same
run that produces the results (Sec.~\ref{sec:ablation}).
